\begin{table}[H]
  \centering
  \footnotesize
  \begin{tabular}{|p{3.6cm}|p{5.0cm}|p{5.0cm}|}
    \hline
    \textbf{Critères} & \textbf{Procédures} & \textbf{Hypothèses} \\
    \hline
    \textcolor{red}{4.1.1} Prise des données & Évaluer la capacité du moyen de stockage. & \\
    \hline
    \textcolor{red}{4.1.2} Interaction à distance fonctionnelle & Évaluer la portée de la transmission et l’accessibilité depuis la centrale. & \\
    \hline
    \textcolor{red}{4.1.3} Détection automatique des mouvements & Évaluer le temps écoulé entre chaque capture. & Un temps d’environ 2 heures et plus entre chaque capture indique que la détection a un problème. \\
    \hline
    \textcolor{red}{4.1.4} Autonomie de la source d’énergie & Évaluer le temps d’autonomie du système. & \\
    \hline
    \textcolor{red}{4.2.1} Installation du système & Calculer le volume et la masse du système. & Un système plus compact et léger est plus facile à installer. \\
    \hline
    \textcolor{red}{4.2.2} Risques de dysfonctionnement & Évaluer la résistance à des conditions extrêmes. & \\
    \hline
  \end{tabular}
  \caption{Plan d’étude pour les différents critères}
  \label{tab:criteres-detail}
\end{table}

\begin{table}[H]
  \centering
  \footnotesize
  \begin{tabular}{|p{3.6cm}|p{5.0cm}|p{5.0cm}|}
    \hline
    \textbf{Critères} & \textbf{Procédures} & \textbf{Hypothèses} \\
    \hline
    \textcolor{red}{4.3.1} Résolution et débit d’images de la caméra & Calculer le nombre d’images par seconde et le nombre de pixels. & \\
    \hline
    \textcolor{red}{4.3.2} Éclairage adéquat de la boîte & Évaluer l’intensité lumineuse présente dans la boîte. & \\
    \hline
    \textcolor{red}{4.3.3} Champ de vision de la caméra & Évaluer le champ de vision selon la longueur focale et la hauteur du capteur. & \\
    \hline
    \textcolor{red}{4.4.1} Coûts de l’équipement & Calculer la somme des coûts des composants du système. & \\
    \hline
    \textcolor{red}{4.4.2} Coûts de déplacement et d’installation & Évaluer les coûts en fonction de la facilité de transport du système. & \\
    \hline
    \textcolor{red}{4.5.1} EDI & Évaluer selon trois dimensions : équité, diversité et inclusion. & Par exemple, un système pouvant mettre en péril la vie de la faune arctique ne serait pas inclusif. \\
    \hline
    \textcolor{red}{4.5.2} Mesure passive (sans pièges) & Évaluer le nombre de perturbations sur la faune. & \\
    \hline
  \end{tabular}
\end{table}
